\section{Returning an automorphic ordinary point to the original ideal}
\label{sec:hida-point-return-v45}

This chapter proves the original \texttt{PointTheoremBaseObject}.
The final implication is restricted to ranks $2,3,4$, exactly as in
the source.  Its ordinary return construction is valid in every
positive rank.  All coefficient maps, the fixed algebraic isomorphism
$\iota:\overline{\mathbb Q}_\ell\simeq\mathbb C$, the integral type,
the auxiliary depth-one subgroup, and the original ideal are retained.
No containment in the kernel of the point is an input to the return.

Fix the original $\mathrm{HypothesesH}'$, Hida setting $s$, level $V$,
auxiliary data $Q$, coefficient DVR $O=c.O$, residual datum $r$,
universal ordinary determinant ring $B$, weight ring $\Lambda$,
unitary Hecke datum $R_H$, and $R_H.\mathrm{IsJ}(B,K,J)$.
In the return and final theorem assume $V.S'=s.S$.
Write $p=V.\mathrm{pr}:G_L\to G_0$ and $\varepsilon=\mathrm{cycBar}$.
The original setting assumes $\ell$ odd; $\ell=3$ is included.
Every local correspondence mentioned below means every map satisfying
the full original $\mathrm{IsRecT}$, with its actual local Artin map.
OFL1 identifies that map with the stored map $a$.

Here are the normalizations used throughout.  At a selected place $w$,
put $q=\#k_w$, choose the stored uniformizer $\varpi$, and set
\[
 \begin{split}
 \mathrm{PS}(\omega)&=\operatorname{Ind}_{B_w}^{\mathrm{GL}_n(L_w),\,\mathrm{norm}}
       \bigl(\omega_i|\,\cdot\,|^{(n-1)/2}\bigr)_{i=0}^{n-1},\\
 \mathrm{psEig}(\omega,j)&=\prod_{0\leq i<j}\omega_i(\varpi)q^{-i},
 \qquad \mathrm{lowW}(\lambda,i)_\tau=\lambda_\tau(n-1-i),\\
 h_i(\tau)&=\mathrm{htW}(\lambda,i)_\tau
                       =\mathrm{lowW}(\lambda,i)_\tau+i,\\
 \mathrm{lowVal}(\lambda,w,j)&=
       \iota\left(\prod_{i<j}\mathrm{algPow}_w
                            (\mathrm{lowW}(\lambda,i))(\varpi)\right).
 \end{split}                                                   \tag{HPT.1}
\]
Thus $h_i(\tau)$ is strictly increasing in $i$.
Our Artin convention is $\operatorname{ord}_w a(s)=-\deg s$;
on the Weil group
\[
 \varepsilon(s)=q^{\deg s}\mathrm{algPow}_w(1)(a(s)).              \tag{HPT.2}
\]
The predicate $\mathrm{OrdChar}(\omega,h,\psi)$ means smoothness of
$\omega$, continuity of $\psi$, and, for every Weil element,
\[
 \psi(s)=\iota^{-1}\omega(a(s))\,
                                \mathrm{algPow}_w(-h)(a(s)).    \tag{HPT.3}
\]
All values in (HPT.3) use the specified embedding and the actual
inclusion of the Weil group into the decomposition group.

\paragraph{Retained foundations.}
We retain exactly the original local classification, support,
local Artin reciprocity, period and local-monodromy foundations
already recorded in GUA, OFL and ITA; the existence of a full
$\mathrm{IsRecT}$ map is retained, as there.  We also retain
$\mathrm{HypothesesH}'.\mathrm{descU}:\mathrm{DescWeak}$ and the
original Jacquet--Shalika isobaric rigidity input
$\mathrm{litN.aut.js}.1$.  Neither is declared newly proved here.
The former applies only to $2\leq n\leq4$, a CM field unramified
over its maximal real subfield at every finite place, and a finite
exception set consisting of split places and containing all places
above $\ell$.  It returns a unitary-group automorphic representation
unramified at nonsplit places and a \emph{cofinite split-place}
Satake comparison.  HPT9 proves the stronger comparison needed here
from that precise output.  The latter input identifies two finite
positive-rank cuspidal families when their unions of normalized
Satake parameters agree at cofinitely all finite places.

\subsection{Point uniqueness and local reverse comparison}

% BEGIN HPT1
\paragraph{HPT1. The evaluated pair determines the entire point.}
\label{claim:HPT1}
Two $O$-algebra points $x,y:B\to\overline{\mathbb Q}_\ell$ are equal
if their evaluated determinant polynomial laws and all their ordered
characters $\psi_{v,i}$ agree.
\begin{proof}
ASP2 puts both entire images in the integer ring $O_E$ of one finite
extension $E/\mathbb Q_\ell$.  It also proves that maximal-ideal
elements have norm less than one.  Both residue maps therefore have
image in the prescribed copy of $k_O$.  Set
\[
                A_0=\{a\in O_E:\bar a\in k_O\}=O+\mathfrak m_E.
\]
This is a finite $O$-module, being an $O$-submodule of the finite
module $O_E$.  Its units are exactly the elements outside
$\mathfrak m_E$: the inverse of a nonzero residue in $k_O$ again
has residue in $k_O$.  It is consequently local with maximal ideal
$\mathfrak m_E$ and residue field $k_O$.  It is Noetherian and
complete: finite modules over the complete DVR $O$ are complete,
and the $\mathfrak m_{A_0}$-adic and $\mathfrak m_O$-adic topologies
agree.  For the last assertion, $A_0/\mathfrak m_OA_0$ is an
Artinian local ring, so a power of its maximal ideal is zero.
Thus $A_0$ satisfies the literal original $\mathrm{IsCO}(O,A_0)$.

Regard $x,y$ as local maps to $A_0$.  Mapping $B.u$ along either
map gives a pair over $A_0$: continuity follows because a local
map sends each maximal-ideal power into the corresponding power;
the residual, polarization, ordinary and inertia relations follow
by applying the map.  The two resulting pairs agree.  Indeed the
characters agree in the subring $A_0\subset\overline{\mathbb Q}_\ell$.
For the determinant law, evaluate on a generic finite group-algebra
sum $\sum T_i g_i$.  Its value lies in $A_0[T_1,\ldots,T_m]$,
which injects into $\overline{\mathbb Q}_\ell[T_1,\ldots,T_m]$.
The assumed full-law equality therefore identifies all coefficients
already in $A_0$.  Evaluation on any scalar algebra proves equality
of the laws over $A_0$.  The uniqueness clause of the original
universal property of $B$ now gives $x=y$.  This argument uses no
topological-generation assertion about selected traces.
\end{proof}
% END HPT1

% BEGIN HPT2
\paragraph{HPT2. Trace relations give the original inertia type.}
\label{claim:HPT2}
Suppose $x.\mathrm{IsRep}(\rho)$ and $\rho$ is semisimple.  If a
subgroup is killed by $p$, it is killed by $\rho$.  At each original
type place $u\in s.S_t$, and for \emph{every} original Teichm\"uller
family $t$, the full inertia action of $\rho$ is conjugate to the
ordered diagonal family $t_i\circ p$ after the specified coefficient
map.  Its characters have open kernels.  Every corresponding
away-$\ell$ Weil--Deligne parameter has precisely this inertia type.
\begin{proof}
If $p(g)=1$, the full-law equality gives
$\operatorname{tr}(\rho(h)(\rho(g)-1))=0$ for every $h\in G_L$.
The algebra generated by a semisimple representation over an
algebraically closed field is a product of full matrix algebras on
its distinct simple constituents.  The trace pairing of its given
representation is the sum of the ordinary matrix trace pairings,
with positive integral multiplicities.  It is nondegenerate in
characteristic zero.  Thus $\rho(g)=1$.

The same proof applied to the literal $B.u.I$ relation shows that
$\rho$ kills the common kernel of the residual characters on
inertia at $u$.  The quotient embeds into a product of copies of
$k_O^\times$, hence is finite abelian of exponent dividing
$N=|k_O|-1$.  Its characteristic-zero representation is a sum of
characters of order dividing $N$.  Residual characteristic
polynomials, given by $B.u.\mathrm{red}$, agree with those of the
ordered residual family.  Reduction is injective on $N$-th roots
of unity, since $\ell\nmid N$.  The eigenvalue multisets therefore
agree with the Teichm\"uller lifts at every inertia element.
Linear independence of distinct characters of the finite abelian
quotient identifies their multisets.  Ordering the summands gives
one conjugating matrix valid on \emph{all} inertia elements.

For completeness, each lifted character has open kernel.  The
residual kernel is open.  On it a Teichm\"uller value $z$ satisfies
$z^N=1$ and $z\equiv1$; the factor $1+z+\cdots+z^{N-1}$ is a unit,
so $(z-1)(1+\cdots+z^{N-1})=0$ forces $z=1$.
For an away-$\ell$ $\mathrm{IsWDOf}$ witness, finite inertia makes
monodromy zero.  Its raw Weil action agrees with $\rho$ on inertia.
The Frobenius semisimplification correction commutes with inertia
and is unipotent; finite-order inertia matrices are semisimple,
so this correction does not change them.  This is the literal
inertia conclusion used in GUA14.  Applying $\iota$ gives the
original $\mathrm{IsTypeWD}$ for every allowed correspondence.
\end{proof}
% END HPT2

% BEGIN HPT3
\paragraph{HPT3. Locally algebraic characters give the exact ordinary characters.}
\label{claim:HPT3}
A continuous decomposition-group character $\psi$ is locally
algebraic of weight $h$ precisely when there is a smooth $\omega$
satisfying (HPT.3).  For such $\omega$, the finite inertia part
$\phi$ in $\mathrm{IsWtPart}$ satisfies
\[
                 \iota\phi(\gamma)=\omega(a(\gamma))^{-1}.
                                                               \tag{HPT.4}
\]
Two continuous characters satisfying (HPT.3) with the same
$\omega,h$ are identical on the entire decomposition group.
\begin{proof}
Given the locally algebraic decomposition, on the Weil group put
\[
       \xi(s)=\iota(\psi(s)/\mathrm{algPow}_w(-h)(a(s))).
\]
On inertia its kernel is exactly the open kernel of the finite
part.  Local Artin reciprocity therefore factors this smooth
character through $a$, giving $\omega$ and (HPT.3).
Conversely define $\phi$ by (HPT.4).  Its kernel is open by
smoothness of $\omega$ and continuity of the Artin map on inertia.
Multiplying (HPT.3) by $\phi$ gives exactly the algebraic-weight
identity in $\mathrm{IsWtPart}$.  This identity holds for every
permitted local Artin map, since OFL1 identifies all of them with
the stored map.  These arguments prove both directions with the
same finite part.  Finally the Weil group is dense in the
decomposition group.  The equality in (HPT.3) and continuity of
the two characters extend their equality from that dense subgroup.
No topology is transported through $\iota$.
\end{proof}
% END HPT3

% BEGIN HPT4
\paragraph{HPT4. Strict weights and the forbidden cyclotomic ratio exclude linkage.}
\label{claim:HPT4}
Suppose $h_i(\tau)$ is strictly increasing in $i$, each
$\mathrm{OrdChar}(\omega_i,h_i,\psi_i)$ holds, and
$\psi_i\ne\psi_j\varepsilon$ for every $i,j$.
Then for $i\ne j$ the characters $\omega_i,\omega_j$ are distinct
and neither is the other times $|\,\cdot\,|^{\pm1}$.
\begin{proof}
A continuous character of a compact group into
$\overline{\mathbb Q}_\ell^\times$ has norm one.  Powers of an
element and its inverse lie in a compact set, so both norms must
be at most one.  If on the Weil group such a character is
$\mathrm{algPow}_w(d)\circ a$, choose a Weil element mapping to
the rational element $\ell\in L_w^\times$.  Its norm is
$|\ell|_\ell^{\sum_\tau d(\tau)}$, so
$\sum_\tau d(\tau)=0$.  The local embedding set is finite and
nonempty because $w\mid\ell$.

If $\omega_i=\omega_j$, (HPT.3) makes $\psi_i/\psi_j$ equal to
the algebraic character with exponents $h_j-h_i$.  Strict increase
makes all of these exponents have the same nonzero sign, contradicting
their zero sum.  If $\omega_i=\omega_j|\,\cdot\,|$, equations
(HPT.2)--(HPT.3) make $\psi_i/(\psi_j\varepsilon)$ have exponents
$h_j-h_i-1$.  For $i>j$ they are all negative.  For $i<j$ they are
all nonnegative, so their zero sum forces every one to vanish.
The ratio is then one on the Weil group and, by density, on the
decomposition group.  This is the prohibited ratio.  Interchanging
$i,j$ treats the inverse norm twist.
\end{proof}
% END HPT4

% BEGIN HPT5
\paragraph{HPT5. A given flag determines jumps and Weil polynomials at every bridge.}
\label{claim:HPT5}
Let $\rho_w$ have a full ordered flag with characters $\psi_i$
satisfying $\mathrm{OrdChar}(\omega_i,h_i,\psi_i)$.
Suppose the original $\mathrm{DRat}$ datum for $\rho_w$ has jumps
$H_\tau$ and Weil characteristic polynomial $P_s(X)$.
Then
\[
 H_\tau=\{h_i(\tau):0\leq i<n\},\qquad
 P_s(X)=\prod_{i=0}^{n-1}(X-\omega_i(a(s))).                    \tag{HPT.5}
\]
These identities hold for every original period ring, bridge,
induced embedding and Weil element in the definition of
$\mathrm{DRat}$.
\begin{proof}
Fix any of those original choices.  The supplied $\mathrm{DRat}$
already makes $\rho_w$ de Rham; closure under subobjects and
quotients makes all terms of its given flag de Rham.  We do not
infer this from the characters alone.  The rank-one calculation OFL4
applies to each character, using HPT3 to identify the actual finite
part: its de Rham jump is $h_i(\tau)$ and its raw Weil character is
$\omega_i\circ a$.  The original period foundations give exact
de Rham and potentially semistable functors on these representations,
with strict induced filtrations.  One may choose a common finite
coefficient field and common finite extension making the finitely
many successive terms semistable; their existence and coefficient
compatibility are the same retained laws used in OFL12.
Apply exactness successively to the supplied flag.  Strictness says
that the multiplicity of each jump is the sum of its multiplicities
in the one-dimensional quotients.  In a basis adapted to the
resulting Weil flag the raw action is upper triangular, with these
rank-one diagonal characters, so its characteristic polynomial is
the product in (HPT.5).  Frobenius semisimplification preserves that
polynomial.  Descent of the chosen finite extension preserves the
flag maps and the polynomial identities.  Finally $\mathrm{DRat}$
uses these same functors and the same representation; uniqueness
of their filtered dimensions and characteristic polynomials gives
the displayed identities.  Since the choices at the start were
arbitrary, the entire original quantifier is preserved.
\end{proof}
% END HPT5

% BEGIN HPT6
\paragraph{HPT6. Recovering the principal series from a toral parameter.}
\label{claim:HPT6}
Let $\tau$ be an actual irreducible smooth local representation with
the original $\mathrm{PlainTor}$ property.  Suppose for every
$\mathrm{IsRecT}$ map the parameter of $\mathrm{algNorm}(\tau)$
has characteristic polynomial $\prod_i(X-\omega_i(a(s)))$ at
every Weil element.  Then $\tau$ is a subquotient of
$\mathrm{PS}(\omega)$.  If the characters satisfy the conclusion
of HPT4, $\tau\simeq\mathrm{PS}(\omega)$.
\begin{proof}
Choose a cuspidal support of $\mathrm{algNorm}(\tau)$.
The original $\mathrm{PlainTor}$ statement makes all its members
characters.  The multisegment exhaustion and support laws of the
original local classification, together with the product law of
$\mathrm{IsRecT}$, identify the characteristic polynomial of its
parameter with the product of the rank-one support parameters.
The rank-one correspondence law makes these parameters
$\xi\circ a$.  OFL1 ensures that the Artin map here is precisely
the one in the hypothesis.

Taking traces gives equality of sums of the support characters
with $\sum_i\omega_i$ on $L_w^\times$, since $a$ is onto.
Linear independence of distinct characters identifies their
multisets, including repetitions.  Undoing
$\mathrm{algNorm}=|\det|^{(1-n)/2}\otimes(-)$ shows that the
support of $\tau$ is the multiset
$\{\omega_i|\,\cdot\,|^{(n-1)/2}\}_i$.  The support theorem
therefore places $\tau$ among the irreducible subquotients of
the normalized induction of these characters, which is exactly
our $\mathrm{PS}(\omega)$.

Under HPT4 these singleton segments are pairwise unlinked.
The original totally disjoint induction theorem, starting with
rank one and applying induction in stages, makes their product
irreducible.  An irreducible subquotient of this irreducible
representation is the representation itself.  This proves the
claimed actual intertwining isomorphism, not merely a support
or semisimplified-parameter equality.
\end{proof}
% END HPT6

% BEGIN HPT7
\paragraph{HPT7. The original automorphic window gives plain torality and reverse unramifiedness.}
\label{claim:HPT7}
If $\mathrm{WinAt}(\pi,t,w)$ holds, every actual local component
$\tau$ of $\pi$ is $\mathrm{PlainTor}$ and every
$\mathrm{rc}(\mathrm{algNorm}(\tau))$ is toral with zero
monodromy.  At $w\nmid\ell$, if this is the Weil--Deligne
parameter of a representation trivial on full inertia, then
$\tau$ is spherical.
\begin{proof}
The literal $\mathrm{WinAt}$ first makes every support member a
rank-one character and then puts the norm of its uniformizer value
strictly between $q^{t-1/2}$ and $q^{t+1/2}$.  Two such characters
cannot differ by $|\,\cdot\,|^{\pm1}$: their norm ratio would be
$q$ or $q^{-1}$, contradicting the two strict inequalities.
Equal characters cause no linkage of singleton segments.
The local component of a cusp is generic by BCS5.  Apply the
original $\mathrm{GenericUnlinkedGL}$ classification: it is an
irreducible product of mutually unlinked segments.  Each segment
has rank-one cuspidal support, and a segment of length at least
two would contribute an adjacent pair differing by
$|\,\cdot\,|$, which the strict window forbids.  All segments
therefore have length one, including any repeated identical
singletons.  Thus this actual irreducible product is the full
principal series of the support characters.  This uses the
generic unlinked-segment theorem, rather than the stronger
pairwise-distinct hypothesis in the totally disjoint theorem.
It proves $\mathrm{PlainTor}$
with its actual inducing data.  The rank-one and product laws of
every $\mathrm{IsRecT}$ give a direct sum of rank-one Weil
characters, with zero monodromy.  Commutators act identically,
which is the original torality condition.  GUA2 supplies $\mathrm{Gd}$.

Now suppose the full Galois inertia is trivial.  Its raw
Weil--Deligne inertia is trivial and stays so under Frobenius
semisimplification.  In the just computed direct sum each support
character is consequently trivial on Artin images of inertia.
These images are precisely $O_{L_w}^\times$.  The common
normalizing norm twist is also trivial there, so all inducing
characters are unramified.  In their normalized principal series,
the Iwasawa formula gives a nonzero $\mathrm{GL}_n(O_{L_w})$-fixed
function with value one at the identity.  The irreducible induction
is $\tau$, proving its sphericity.  This proves the reverse
implication rather than appealing to injectivity of a spherical
parameter assignment.
\end{proof}
% END HPT7

% BEGIN HPT8
\paragraph{HPT8. Compatibility for the supplied representation and its supplied flag.}
\label{claim:HPT8}
Suppose $\rho$ is continuous and $\mathrm{IsAutCsdA}(\iota,\rho)$.
Choose the regular algebraic conjugate-self-dual cusp $\pi$ in its
actual $\mathrm{AutCsd0}$ witness.  For every $w\nmid\ell$, and
for every $w\mid\ell$ where the supplied flag has no ratio
$\psi_i=\psi_j\varepsilon$, the actual local components are
$\mathrm{PlainTor}$ and the full $\mathrm{LocGlobT}$ comparison
holds for this same $\rho$.  At a selected $\ell$-place with an
arithmetic ordered flag of weights $h_i$, the component is
$\mathrm{PS}(\omega)$ with (HPT.3), and its Hodge multiset is
$\{h_i\}_i$.
\begin{proof}
The $\mathrm{AutCsd0}$ witness includes semisimplicity and the
cofinite attachment of this $\rho$ to $\pi$.  Polarize $\pi$ by
the original character $\delta_{L/L^+}^{n}$: its pullback along
the norm is one and its sign is $(-1)^n$.  Thus every hypothesis
of the original RowT application GUA7 holds.  It produces a
continuous semisimple representation $\rho_0$, its Hodge
multisets and its full guarded local compatibility.
At cofinitely many places the cusp is spherical and away from
$\ell$.  GUA12 and HED8 with $t=(1-n)/2$ convert the RowT
comparison into its cofinite attachment to $\pi$, with the
original algebraic Satake normalization.  Its Frobenius
polynomials thus equal those of $\rho$.  Chebotarev
and continuity identify the characteristic polynomials everywhere;
Brauer--Nesbitt gives one constant intertwiner
$\rho\simeq\rho_0$.  Enlarge a finite coefficient field to contain
this matrix and its inverse.  Tensoring that matrix with every
period ring transports the filtered modules and Weil actions,
so the full Hodge and local comparison statements hold for $\rho$.
This is the same finite-coefficient descent used in ADK1 and GUA10;
it preserves all bridges and every $\mathrm{IsRecT}$ quantifier.

Continuity of each character in the supplied flag follows by
conjugating the matrix representation into upper triangular form
and taking a diagonal entry.  A flag at one prolongation supplies
the same kind of flag at every prolongation: conjugate the
decomposition subgroup by a global Galois element and conjugate
the flag basis by its matrix.  This group isomorphism is onto and
the cyclotomic character is conjugation invariant.  The forbidden
ratio therefore stays forbidden.  Hence the precise $\mathrm{Unl1}$
premise of $\mathrm{WitA}$ holds.  Its first field upgrades this
to $\mathrm{PotUnl}$; its second field provides a single window
for this attached cusp at all away-$\ell$ places and these
above-$\ell$ places.  HPT7 gives plain torality and $\mathrm{Gd}$,
so RowT's toral, zero-monodromy branch supplies the stated
$\mathrm{LocGlobT}$.  No rank-dependent alternative branch is used.

Arithmeticity gives local algebraicity of the ordered characters.
HPT3 supplies $\omega_i$.  Apply HPT5 to the actual Hodge and
local comparison just obtained, then HPT6 and HPT4.
This gives the full principal series and the ordered-weight
Hodge multiset.  All steps use the original $\rho$, with no
assumption that the point already kills $J$.
\end{proof}
% END HPT8

% BEGIN HPT9
\paragraph{HPT9. Upgrading weak descent to every split place and infinity.}
\label{claim:HPT9}
Let $\pi$ be the cusp of HPT8 and let $\sigma$ be a unitary-group
automorphic representation with the cofinite split Satake
comparison supplied by $\mathrm{DescWeak}$.  Then every actual
split local component of $\sigma$ is the corresponding component
of $\pi$, up to an actual intertwiner, and the infinity multiset
of $\sigma$ equals that of $\pi$.
\begin{proof}
Use BCS18 and GUA18 on this actual $\sigma$.  Write its cuspidal
pieces as $\Pi_b$ of ranks $a_b>0$, with lengths $d_b>0$, and
write $R_\sigma$ for the continuous semisimple representation
produced there.  At the cofinite split places where all objects
are spherical and away from $\ell$, GUA14 and the weak comparison
identify the characteristic polynomials of $R_\sigma$ and
$\rho$.  Every degree-one place is split by HED6.  HED7 therefore
extends this equality to every element of $G_L$.

At cofinitely \emph{all} finite places, BCS19 and GUA14 compute
the same polynomial as the union of the strings
\[
       \{\alpha q_w^{k-(d_b-1)/2}:
                  \alpha\in\mathrm{Sat}(\Pi_b,w),\ 0\leq k<d_b\}.
                                                               \tag{HPT.6}
\]
The attachment of $\rho$ to $\pi$ identifies (HPT.6) with
$\mathrm{Sat}(\pi,w)$.  Interpret the left side as the family
of actual cuspidal twists
$\Pi_b\otimes|\det|^{(d_b-1)/2-k}$ indexed by $(b,k)$.
The Satake sign is correct because twisting by $|\det|^u$
multiplies a Satake parameter by $q_w^{-u}$.
The retained Jacquet--Shalika rigidity input applies to these
two finite positive-rank cusp families over the totally complex
field $L$.  It gives a bijection of this index set with the
singleton indexing $\pi$.  Since all $d_b$ are positive, there
is exactly one $b$, its $d_b$ is one, its rank is $n$, and
its cusp is $\pi$ itself.

For these witnesses the full split product in BCS18 is just the
length-one local component of $\pi$.  The length-one Speh and
one-factor induction maps are identities.  Local component
uniqueness LTP3 identifies every other actual split component
of $\sigma$ with it.  The BCS18 infinity equality becomes
exactly the infinity multiset of $\pi$.  This argument proves
the strong comparison for this weak descent; it does not assert
that arbitrary unitary-group automorphic representations have
only length-one Speh factors.
\end{proof}
% END HPT9

\subsection{Constructing the nonzero vector in the original coefficient space}

% BEGIN HPT10
\paragraph{HPT10. Simultaneous global compact fixed vectors.}
\label{claim:HPT10}
Let $\sigma$ be the actual irreducible automorphic unitary-group
representation.  Assume it is unramified outside $V.S'$ and its
actual component at $v_0$ has a nonzero $K_c(1)$-fixed vector.
There is a nonzero vector fixed by the original subgroup $K_0$
which is maximal compact outside $V.S'$, is $K_c(1)$ at $v_0$,
and is the identity at the selected $\ell$-places, type places
and auxiliary places.
\begin{proof}
For a finite place of $L^+$, restriction to its coordinate factors
and replacement of the other coordinates by one gives a retraction
of the finite adelic unitary group onto its local factor $H_w$.
The unitary equation is preserved because the entire pair of
conjugate coordinates is retained.  Thus the group is a product
of $H_w$ with its commuting complementary factor.  The retraction
preserves $K_0$.  Its local factors in $K_0$ are compact: the
integral unitary factor is a closed subgroup of the corresponding
matrix compact group, and $K_c(1)$ is compact as well.

We first record a fixed-vector transfer valid also at nonsplit
places.  If $W$ is a nonzero $H_w$-stable subspace of $\sigma$,
global irreducibility says that the span of its global translates
is $\sigma$.  By the product decomposition it is the sum of
translates $cW$ by the complementary factor.  Normalized averaging
over a compact local subgroup $C_w$ is an exact projector on a
smooth complex representation: on each vector its orbit map is
locally constant, so its integral is a finite sum.  This projector
commutes with $c$.  If $W^{C_w}=0$, the projector vanishes on every
$cW$ and hence on $\sigma$.  Therefore $\sigma^{C_w}\ne0$
implies $W^{C_w}\ne0$.  At the places under consideration the
required nonzero fixed vector in $\sigma$ comes from the stated
unramified condition or from the actual $v_0$ component; LTP5
transfers the latter vector to $\sigma$.

Choose any nonzero smooth vector in $\sigma$.  Its open stabilizer
and the restricted-product topology give a finite set $T$ of
places such that every element of $K_0$ which is one at $T$ fixes
it.  Induct on $|T|$.  For $w\in T$, the nonzero subspace fixed
by these elements of $K_0$ is $H_w$-stable, because the supports
are disjoint.  The preceding averaging argument produces in it
a nonzero vector fixed also by $K_0\cap H_w$.  Every element of
$K_0$ which is one at $T\setminus\{w\}$ is a product of its
local retraction at $w$ and an element which is one at all of $T$.
The new vector is fixed by both.  Removing one place at a time
gives the asserted nonzero $K_0$-fixed vector.  This proof requires
no choice of a nonsplit local Langlands parameter.
\end{proof}
% END HPT10

% BEGIN HPT11
\paragraph{HPT11. The integral type and the infinity dual give the actual nonzero part.}
\label{claim:HPT11}
Assume the hypotheses of HPT10, the original infinity weights
$\lambda_\tau^\vee$, and, at every type place, an actual local
component whose parameter has the prescribed type for every
$\mathrm{IsRecT}$ and every original Teichm\"uller family.
For the original $\mathrm{TypeForAlt}$ representation $\tau_0$,
every specified algebraic basis $e$ with $\mathrm{IsBasis}(\iota,e,\xi)$
gives
\[
                         Y\mathrm{Of}(Q,\iota,\tau_0,e,\xi,\sigma)\ne0.
                                                               \tag{HPT.7}
\]
\begin{proof}
Choose the original Teichm\"uller family using ITA2.  Put
$W=\sigma^{K_0}$, which is nonzero by HPT10.  The actual product
section $J_t$ of the type groups supplied by LTP16 commutes with
$K_0$.  Local component isotypy LTP3 and the given local parameter
types permit the full product type-detection statement ITA12 on
$W$.  It gives a nonzero column $F_0=(F_{0,i})\in W^m$ with
\[
       \sum_a\tau_C(k)_{ia}\,R_U(J_tk)F_{0,a}=F_{0,i}
                   \quad(k\in K_T, 0\leq i<m),                \tag{HPT.8}
\]
where $\tau_C$ is the specified coefficient map applied to
$\tau_0$.  The inverse transpose in ITA9 is exactly what produces
this column equation; no dual of $\tau_0$ is substituted.

Let $N\subset W^m$ be the space of all columns satisfying
(HPT.8).  It is nonzero and stable under the infinity group,
since finite and infinite adelic factors commute.  Every selected
infinity factor on each of its components has the given dual
highest weight.  The source embeddings index the infinity places
bijectively.  LTP13 proves injectivity.  For surjectivity, at each
selected split place $w\mid v\mid\ell$ their number is
$[L_w:\mathbb Q_\ell]=[L_v^+:\mathbb Q_\ell]$: characteristic
zero gives exactly that many embeddings of the local field,
all continuous for its unique extended absolute value.
Summing over $v\mid\ell$ gives $[L^+:\mathbb Q]$, by the
finite algebra decomposition of $L^+\otimes\mathbb Q_\ell$.
This is also the number of infinite places of the CM field $L$.
The injective map between these finite sets is therefore onto.
LTP13's explicit dual identification,
and the successive-Hom argument of ITA11, therefore give a
nonzero intertwiner
\[
       \Phi:D^*\longrightarrow N,\qquad
       D=\bigotimes_\tau\nabla_{\lambda_\tau},\qquad
       R_U(z)\Phi(\ell)=\Phi(\ell\circ\xi_\infty(z^{-1})).       \tag{HPT.9}
\]
More explicitly, take a nonzero copy of the required irreducible
dual at the first factor; its Hom multiplicity space remains
isotypic for every commuting remaining factor.  Repeating over
the finite set and evaluating the tensor of these Hom maps gives
(HPT.9).  At each step a nonzero map out of an irreducible module
is injective, so the final map is nonzero.  This also proves that
no finite-factor condition was lost.

Apply the actual inverse scalarization ADP8--ADP9 to each
component of $\Phi$, using the specified basis $e$.  It gives
coefficient-valued functions $F_i$ in the original $\mathrm{AF}(\xi)$
such that every scalarization of $F_i$ is $\Phi(\ell)_i$.
It remains to check the \emph{actual} $U_Q$ relation.  For
$u\in U_Q$, lift $((u: G^a),1)$ to a finite adelic element $z$
using the surjectivity in ADP9.  Coordinate by coordinate,
\[
                      z=J_t(\mathrm{toKT}(u))\,k_0,
                                      \qquad k_0\in K_0.        \tag{HPT.10}
\]
Indeed the selected $\ell$-coordinates of $z$ are one; at type
places the section removes precisely its type coordinate;
outside $V.S'$ the remaining coordinate is integral; at $v_0$
it lies in the original $K_c(1)$; and at auxiliary coordinates
the $U_Q$ condition makes it one.  These are exactly the
defining coordinates of $K_0$.

The translation identity ADP7, the trivial infinity coordinate
of $z$, equation (HPT.8), and the $K_0$ invariance give, after
every scalarization,
\[
       \sum_a\tau_C(\mathrm{toKT}(u))_{ia}
                     F_a(g((u: G^a),1))=F_i(g).               \tag{HPT.11}
\]
Scalarization is injective by ADP7, so (HPT.11) is an equality
of the original coefficient functions.  All their scalarizations
belong to $\sigma$.  Thus $F\in Y\mathrm{Of}$ by its literal
definition.  If $F=0$, all scalarizations and hence $\Phi$ would
vanish, contradicting its construction.  This proves (HPT.7).
\end{proof}
% END HPT11

% BEGIN HPT12
\paragraph{HPT12. A joint eigenfunctional produces a joint eigenvector.}
\label{claim:HPT12}
Let a commuting family $(T_i)$ act on a finite-dimensional vector
space over an algebraically closed field.  If a nonzero functional
$L$ satisfies $LT_i=\mu_iL$ for every $i$, there is a nonzero vector
$v$ with $T_iv=\mu_iv$ for every $i$.
\begin{proof}
The algebra generated by the family is finite-dimensional and
commutative, even if the indexing set is infinite.  Decompose the
space into its finitely many simultaneous generalized-character
spaces.  On a space of character $\chi\ne\mu$, choose $i$ with
$\chi(T_i)\ne\mu_i$.  A power of $T_i-\chi(T_i)$ vanishes on
that space, while applying $L$ multiplies by the nonzero scalar
$(\mu_i-\chi(T_i))$ to the same power.  Hence $L$ vanishes there.
The $\mu$ space is consequently nonzero.  On it the commuting
operators $T_i-\mu_i$ generate the nilpotent maximal ideal of a
finite local algebra.  The last nonzero term of its powers applied
to the space is killed by that ideal.  Any nonzero vector there
is the required simultaneous eigenvector.
\end{proof}
% END HPT12

% BEGIN HPT13
\paragraph{HPT13. A principal series has the exact local ordinary eigenline.}
\label{claim:HPT13}
Suppose each smooth $\omega_i$ is trivial on $U_c(c_0+1)$.
In $\mathrm{PS}(\omega)$ there is a nonzero $I_1(c_0)$-fixed
vector on which the compact diagonal torus acts by
$\prod_i\omega_i(t_i)$ and, for every $0\leq j\leq n$,
the original positive double-coset operator at
$a_j=\operatorname{diag}(\varpi I_j,I_{n-j})$ acts by
$\mathrm{psEig}(\omega,j)$.
\begin{proof}
Realize the principal series by functions $f(bg)=\chi(b)f(g)$,
where
\[
       \chi(t)=\delta_B(t)^{1/2}\prod_i\omega_i(t_i)|t_i|^{(n-1)/2}.
\]
The function supported on $BI_1(c_0)$ with value one at the
identity is well defined: $\chi$ is trivial on $B\cap I_1(c_0)$,
whose diagonal entries lie in $U_c(c_0+1)$.  Thus evaluation at
the identity is a nonzero functional on the $I_1(c_0)$-fixed
space.  That space is finite-dimensional.  Iwasawa decomposition
reduces its functions to finitely many coordinates on
$K/I_1(c_0)$, and $I_1(c_0)$ is open in the compact group $K$.

The compact torus and all positive operators preserve this fixed
space.  OST5 and OST26 give their actual coset representatives
and their pairwise commutation.  Evaluation at the identity has
torus eigenvalue $\chi(t)=\prod_i\omega_i(t_i)$ and positive
operator eigenvalue $[N:aNa^{-1}]\chi(a)$: all representatives
lie in $N$, on which the inducing character is one.
For $a=a_j$, the index and character are
\[
 [N:a_jNa_j^{-1}]=q^{j(n-j)},\qquad
 \chi(a_j)=\prod_{i<j}\omega_i(\varpi)q^{-(n-1-i)}.
\]
Since $\sum_{i<j}(n-1-i)=j(n-j)+\sum_{i<j}i$, their product is
exactly $\mathrm{psEig}(\omega,j)$.  Apply HPT12 to obtain a
simultaneous eigenvector, including all compact-torus operators.
In the source $\mathrm{IsInvOn}(\phi,\omega)$ notation this
torus character is $\phi^{-1}$.  The depth is $c_0+1$ as stated;
no increase in depth or restriction on the residue prime is made.
\end{proof}
% END HPT13

% BEGIN HPT14
\paragraph{HPT14. Simultaneous local eigenlines give an actual global ordinary vector.}
\label{claim:HPT14}
Let $H$ be the original Hecke setting, let $\mathrm{IsSpec}$ be
its actual specialization at an arithmetic point $x_A:A\to
\overline{\mathbb Q}_\ell$, with data $\lambda,\xi,c_0,\mathrm{sp}$,
and let $Y$ be an original nonzero $\mathrm{PartAltAlt}$.
Suppose its selected local components are $\mathrm{PS}(\omega_v)$,
the auxiliary components have nonzero $K_q$-fixed vectors, and
\[
 \omega_{v,i}(u)\,\iota\mathrm{finPart}_{v,i}(u)=1,\qquad
 \left|\iota^{-1}\bigl(\mathrm{lowVal}(v,j+1)^{-1}
                 \mathrm{psEig}(\omega_v,j+1)\bigr)\right|_\ell=1.
                                                               \tag{HPT.12}
\]
Then $Y$ contains a nonzero vector in the actual specialized
ordinary space at $K,c_0$, with all spherical eigenvalues $Y.a$
and all normalized ordinary eigenvalues
\[
             \alpha_{v,j}=\mathrm{lowVal}(v,j+1)^{-1}
                                     \mathrm{psEig}(\omega_v,j+1).
                                                               \tag{HPT.13}
\]
\begin{proof}
The triviality field of $\mathrm{IsSpec}$ and (HPT.12) make each
$\omega_{v,i}$ trivial on $U_c(c_0+1)$.  Apply HPT13 locally.
To impose these conditions simultaneously, start with $Y$ and
process the finite set of selected and auxiliary factors.  The
current subspace is stable and isotypic for every remaining factor:
actions at different selected coordinates commute; selected and
auxiliary actions commute by the literal right-translation
formula; two auxiliary actions commute by the original disjoint
coordinate condition in $H.\mathrm{Ax}$.  Isotypy means that any
nonzero stable subspace contains an embedded copy of its given
irreducible local representation (LTP3).  Map the nonzero vector
of HPT13, or the supplied auxiliary fixed vector, into this copy.
Intersecting with that factor's conditions therefore leaves a
nonzero subspace and preserves all previous conditions.  Finite
induction gives one nonzero $F$ satisfying every condition.

The local $N$ and lower-unipotent fixed conditions combine into
the corresponding product conditions.  The $U_Q$ relation already
holds on $Y$, and the auxiliary fixed conditions extend it to
$H.\mathrm{lev}(K)$.  The torus character in HPT13 and (HPT.12)
is exactly the finite-part character in $H.\mathrm{ClAt}$.
Thus $F\in H.\mathrm{ClAt}(K,c_0)$.  OST26--OST27 compute its
positive Hecke eigenvalues and give (HPT.13).

For ordinary projection one must use the original indexing:
the element $z$ is the product of $a_j$ for $j=0,\ldots,n-1$,
whereas the displayed generators are indexed by $j+1=1,\ldots,n$.
The normalized $z$ eigenvalue on $F$ is
\[
              \prod_v\prod_{j=0}^{n-1}
                  \mathrm{lowVal}(v,j)^{-1}\mathrm{psEig}(\omega_v,j).
\]
The $j=0$ factor is one and all remaining factors have norm one
by (HPT.12).  Hence $F$ belongs to the original unit-eigenvalue
ordinary part.  Its spherical eigenvalues are those of $Y$
because $F\in Y$.  This proves the complete simultaneous assertion
on the original function space.
\end{proof}
% END HPT14

% BEGIN HPT15
\paragraph{HPT15. Exact unit and residue normalization at the original point.}
\label{claim:HPT15}
Let $x:B\to\overline{\mathbb Q}_\ell$ be arithmetic of weight
$h=\mathrm{htW}(\lambda)$, and suppose (HPT.3) holds for its
ordered characters.  Then the first identity in (HPT.12) holds,
and every $\iota^{-1}\alpha_{v,j}$ is a unit value of $x$,
congruent to the specified ordinary residual eigenvalue
$R_H.a(v,j)$ modulo elements of norm less than one.
If $Y$ has the same Frobenius polynomials as $\rho$ with
$x.\mathrm{IsRep}(\rho)$ at the original spherical places, its
spherical eigenvalues have the analogous residual congruence.
\begin{proof}
For a local unit $u$, choose a Weil element mapping under $a$
to $u$.  Its degree is zero.  Insert (HPT.3) in the original
$\theta$ relation and use the exact weight map
$\Lambda.\mathrm{toPair}(B.u)$; (HPT.2) cancels the $i$ copies
of the cyclotomic algebraic factor against
$h_i=\mathrm{lowW}_i+i$.  The remaining equality is
$\omega_{v,i}(u)\iota\mathrm{finPart}_{v,i}(u)=1$.
This calculation is HEP4--HEP6 in the reverse direction and uses
the given finite part, not a new choice of it.

Choose a Weil lift $s_\varpi$ with $a(s_\varpi)=\varpi$ and
degree $-1$, and let $g$ be its actual decomposition-group image.
The unit of $B$ relevant to the $(j+1)$-st operator is
\[
             b_{v,j}=\prod_{i=0}^{j}
                          B.u.\psi_{v,i}(g)\,\varepsilon(g)^i.
\]
Using (HPT.2)--(HPT.3) and (HPT.1) term by term gives
\[
                          x(b_{v,j})=\iota^{-1}\alpha_{v,j}.
                                                               \tag{HPT.14}
\]
Indeed the algebraic exponents are $-h_i+i=-\mathrm{lowW}_i$
and the numerical factors are $q^{-i}$.  ASP2 says that the value
of a unit and its inverse both have norm at most one, so this
value has norm one.  The residual character field of $B.u$ and
the literal $R_H.\mathrm{res.unif}$ identify its residue with
$R_H.a(v,j)$.  ASP2 then gives the strict norm inequality.

At a spherical place $w\nmid\ell$, compare the coefficient of
$X^{n-j}$ in the common Frobenius polynomial.  With the source
Hecke convention it is
$(-1)^j q_w^{j(j-1)/2}Y.a(w,j)$.  The residual determinant law
and $R_H.\mathrm{res.frob}$ give the same coefficient with
$R_H.a(w,j-1)$ for $1\leq j\leq n$.  The factor
$q_w^{j(j-1)/2}$ has $\ell$-adic norm one; division preserves
the strict congruence.  This proves every residual congruence
needed for the actual residual block.
\end{proof}
% END HPT15

% BEGIN HPT16
\paragraph{HPT16. The actual ordinary vector gives a point of the original residual Hecke algebra.}
\label{claim:HPT16}
In HPT14 suppose all its eigenvalues are congruent to the prescribed
residual function $a$.  Then the original residual module
$H.MK(\theta,a,K)$ is nonzero, and there is an $O$-algebra point
$y:H.TK(\theta,a,K)\to\overline{\mathbb Q}_\ell$ extending $x_A$
and taking every generator to the eigenvalue in HPT14, transported
by $\iota^{-1}$.  This is the entire literal
\texttt{coreUp\_alt} conclusion.
\begin{proof}
Let $N$ be the original ordinary module before taking the residual
block.  OHA4 and OHA8 make $N$ finite free and its given operators
commuting.  The span and independence fields of ASP13 identify
$N\otimes_{A,\iota x_A}\mathbb C$ with the specialized ordinary
space.  The nonzero vector $F$ of HPT14 therefore determines an
eigenline in this scalar extension of the actual $N$.
Its character on the actual algebra $T\subset\operatorname{End}_A(N)$
exists because every generator preserves this line.  It is
well defined on $T$: any relation acting as zero on $N$ acts
as zero after scalar extension and hence on that line.
Transport its values through $\iota^{-1}$ to obtain $\chi$.
It extends $x_A$ and takes the generators to precisely the stated
eigenvalues.  All its values are integral: $T$ is generated over
$A$ by these integral eigenvalues, and ASP2 bounds $x_A(A)$ by one.
Reduction of $\chi$ kills $\mathfrak m_A$ and every $T_i-a_i$.

OHA5--OHA6 decompose $T$ into its actual finite collection of
local idempotent factors.  The ideal
$(\mathfrak m_A,T_i-a_i)_i$ singles out at most one factor; the
reduction of $\chi$ shows that this factor exists.  Its idempotent
$e$ maps to one under $\chi$, whereas all other idempotents map
to zero.  Consequently $eF=F\ne0$ and $eN=H.MK(\theta,a,K)$
is nonzero.  The action of $eT$ on $eN$ is faithful: it is zero
on $(1-e)N$ already, so an element killing $eN$ kills $N$ and
is zero in the original operator algebra.  Thus $eT$ is exactly
the original block algebra $H.TK$.  Restrict $\chi$ to it to get
$y$.  This proves the coefficient-map identity and all spherical
and normalized ordinary generator identities simultaneously.
\end{proof}
% END HPT16

% BEGIN HPT17
\paragraph{HPT17. The full original ordinary return theorem in every positive rank.}
\label{claim:HPT17}
Retain the original data with $n\geq1$ and $V.S'=s.S$.
Suppose $x$ is arithmetic with $h=\mathrm{htW}(\lambda)$,
$\rho$ is continuous semisimple and $x.\mathrm{IsRep}(\rho)$.
Suppose an actual automorphic $\sigma$ satisfies all of the
following original return inputs: $\mathrm{SphEq}(\iota,\sigma,\rho)$;
infinity weights $\lambda^\vee$; unramifiedness outside $s.S$;
a nonzero $K_c(1)$-fixed vector in its actual $v_0$ component;
at each selected $\ell$-place an actual component
$\mathrm{PS}(\omega_v)$ satisfying (HPT.3) for $x.\psi_{v,i}$;
and $\mathrm{StrongEquivalenceHida}$ at all type places, with
its quantifier over every actual local component and every
$\mathrm{IsRecT}$.  Then $J\subseteq\ker x$.
\begin{proof}
AFS8 translates the arithmetic weight to the original admitted
$\lambda$.  Use ASP13 to choose the actual specialization map,
algebraic basis, depth and coefficient structure.  HPT2 and the
strong comparison supply all parameter types at the type places,
for every Teichm\"uller family.  HPT11 constructs a nonzero
$Y\mathrm{Of}$ with this basis.  UPA16 now produces the original
part $Y$ and its continuous semisimple frame representation $R_Y$,
with all local components still those of this same $\sigma$.

At each original spherical place $w\in H.W_p$, the hypotheses
make $\sigma$ spherical.  Choose an actual local component and
its Satake parameters.  $\mathrm{SphEq}$ computes the Frobenius
polynomial of $\rho$ from those parameters; the local comparison
of $Y$ and GUA14 compute the polynomial of $R_Y\circ p$ from
the same parameters.  They agree for every original prolongation
and Frobenius witness.  HEP1--HEP2, using the given full determinant
law of $x$ and the original density of $H.W_p$, give the full
frame equality
\[
                              x.D=\det R_Y.                    \tag{HPT.15}
\]
This is an equality of polynomial laws, not only of Frobenius
characteristic polynomials.  HPT15 gives every spherical residual
congruence and every normalized ordinary unit and congruence.
Actual local-component uniqueness identifies $Y$'s selected
components with the given principal series.

There are no auxiliary factors here.  Indeed $Q.\mathrm{range}$
would put the place of any $q\in Q$ in $V.S'=s.S$, contrary
to $Q.\mathrm{notS}$.  Thus the auxiliary fixed-vector requirement
of HPT14 is vacuous.  Apply HPT14--HPT16 to get the nonzero
original residual module and its point $y$ with the exact
weight and eigenvalues.

In the zero-module branch of $\mathrm{IsJ}$ this nonzero module
is a contradiction.  In the other branch take its actual
surjection $f:B\twoheadrightarrow H.TK$ with $J=\ker f$ and
the stated determinant, ordinary-character and weight comparisons.
Apply HEP7 to a nonzero vector of $Y$, the just constructed
Hecke point, and its spherical and ordinary eigenvalue identities.
It gives
\[
                (yf).D=\det R_Y,\qquad
       \mathrm{OrdChar}(\omega_{v,i},h_i,(yf).\psi_{v,i}).       \tag{HPT.16}
\]
The finite-part premise in that application is exactly HPT15,
since $y$ has the weight $x|_\Lambda$.  Equations
(HPT.15)--(HPT.16) identify the two full laws.  HPT3 identifies
all ordered characters.  HPT1 therefore gives $x=yf$ as maps on
the entire original ring $B$.  For $b\in J=\ker f$ we obtain
$x(b)=y(f(b))=0$, proving the conclusion.  This is the literal
\texttt{upU\_w\_alt} endpoint with all of its source inputs,
proved here on the actual module and actual ideal.
\end{proof}
% END HPT17

% BEGIN HPT18
\paragraph{HPT18. The complete low-rank point theorem.}
\label{claim:HPT18}
For every original point $x$, weight $h$ and representation
$\rho:G_L\to\mathrm{GL}_n(\overline{\mathbb Q}_\ell)$ with
$2\leq n\leq4$, assume
\[
 \begin{gathered}
 x.\mathrm{Arith}(h),\qquad x.\mathrm{IsRep}(\rho),\qquad
 \mathrm{Continuous}(\rho),\qquad \mathrm{IsAutCsdA}(\iota,\rho),\\
 \rho|_{I_{s.A(v_0)}}=1,\qquad
 \mathrm{HasFlag}(\rho|_{D_{s.A(v)}},x.\psi_v)
                  \quad\text{for every selected }v\mid\ell,\\
 x.\psi_{v,i}\ne x.\psi_{v,j}\varepsilon|_{D_{s.A(v)}}
                          \quad\text{for every }v,i,j .
 \end{gathered}                                                 \tag{HPT.17}
\]
Then $J\subseteq\ker x$.  Hence the same original $J$ satisfies
the entire \texttt{PointTheoremBaseObject} field.
\begin{proof}
Use AFS8 to obtain the admitted dominant $\lambda$ with
$h_i=\lambda(n-1-i)+i$.  The $\mathrm{IsAutCsdA}$ witness
supplies a regular algebraic conjugate-self-dual cusp $\pi$,
the attachment and semisimplicity of $\rho$, and its $\mathrm{WitA}$
guard.  HPT8 gives the full compatible local parameters and
the selected ordinary principal series for this supplied
$\rho$.  Outside $s.S=V.S'$, the map $p$ kills inertia.
HPT2 makes $\rho$ unramified there; these places are away from
$\ell$.  HPT7 makes every local component of $\pi$ there
spherical, and the original Satake theorem supplies its parameters.

Apply the retained $\mathrm{DescWeak}$ with $E=L$, prime $\ell$
and exception set $s.S$.  The original Hida setting supplies
finite unramifiedness of $L/L^+$, finiteness of $s.S$, splitting
of every place above $s.S$, and inclusion of every place above
$\ell$.  The rank inequalities are exactly (HPT.17).  The
cusp and its Satake condition were just verified.  Thus it gives
an actual $\sigma$, unramified at all nonsplit places, with the
cofinite split comparison.  HPT9 upgrades this to every split
local component and the infinity multiset.

We verify each input of HPT17.  At selected $\ell$-places the
local component is the principal series already obtained in HPT8,
with the same ordered $\omega_i$, and the corresponding Hodge
multiset is $\{\lambda(n-1-i)+i\}_i$.  If $\nu$ denotes the
dominant infinity weight of $\sigma$, the original Hodge/infinity
normalization and HPT9 say that the same multiset is
$\{-\nu_i+i\}_i$.  Both sequences are strictly increasing.
Their minima agree; remove that common term and repeat to get
termwise equality.  Hence $\nu_i=-\lambda(n-1-i)$, exactly
$\lambda^\vee$.  This holds at every source embedding, including
its original local embedding identification.

Outside $s.S$, nonsplit unramifiedness is the descent output.
At split places it follows from HPT9, the spherical components
of $\pi$, and LTP20--LTP21.  At $v_0$ the residue characteristic
is different from $\ell$ by the original Hida-setting condition.
The full-inertia hypothesis in (HPT.17) and HPT7 make its
$\pi$ component spherical.  Its nonzero maximal-compact fixed
vector is fixed by the original subgroup $K_c(1)$.  HPT9
transports this vector to the actual $\sigma$ component.

At every split away-$\ell$ place, the full parameter comparison
of HPT8 and the same-component statement of HPT9 give the
spherical Frobenius law GUA14, which is precisely
$\mathrm{SphEq}$.  At every type place, all actual components
match by HPT9 and all correspondences match by HPT8, giving
the literal $\mathrm{StrongEquivalenceHida}$, with no restriction
to one correspondence or one choice of component.  All hypotheses
of HPT17 have now been proved from (HPT.17) and the original
data.  Its conclusion is $J\subseteq\ker x$, as required.
In particular this containment has not been used in producing
$\sigma$, the local vectors, or the point $y$.
\end{proof}
% END HPT18
